% ==========================================================================
% Planning reference: what each thesis chapter and section should contain.
% Heading levels match the actual thesis for easy cross-referencing.
% ==========================================================================

%--- ~5 pages ---%
\chapter{Introduction}\label{ch:introduction}
\markboth{Introduction}{}

\section{Motivation and Context}
\begin{itemize}
    \item Industrial textile recycling requires automated sorting by garment condition (intact, minor damage, major damage, unusable)
    \item Current approaches rely on manual inspection or single-robot rigid-object pipelines; neither scales to deformable textiles on conveyors
    \item Multi-agent robotic systems combining pick, stretch, inspect, and sort offer a path to full automation
    \item This thesis builds on the simulation environment and validated tendon-driven robot from the preceding project thesis (Meyer, 2026)
    \item Central challenge: coordinating two robots under RL for deformable cloth manipulation with vision-based classification
\end{itemize}

\section{Problem Statement and Objectives}
\begin{itemize}
    \item Thesis problem: developing a multi-agent RL system that coordinates a tendon-driven picker and a collaborative sorting arm for condition-based textile sorting in simulation
    \item Research objectives:
    \begin{enumerate}
        \item Selective sorting task: retrieve T-shirts from mixed conveyor waste
        \item Deformable cloth simulation: integrate T-shirt cloth physics (PBD/VBD)
        \item Second robot arm: stretch garments for camera-based inspection
        \item Damage classification: vision-based condition assessment
        \item Multi-agent coordination: joint policies for pick-stretch-classify-sort
        \item Sim-to-real preparation: domain randomization, ROS 2 integration pathway
    \end{enumerate}
    \item Scope: simulation-based development and evaluation; physical deployment is prepared but not executed
\end{itemize}

\section{Thesis Outline}
\begin{itemize}
    \item Brief summary of chapters 2--7 for reader orientation
    \item Explicitly reference the project thesis as the foundation (simulation, robot model, basic RL tasks)
\end{itemize}

%--- ~20--25 pages ---%
\chapter{Background}\label{ch:background}
\markboth{Background}{}

\section{Summary of Foundation (Project Thesis)}
\begin{itemize}
    \item Brief recap of the tendon-driven robot model, simulation environment, and RL pipeline established in the project thesis (Meyer, 2026)
    \item Reference rather than repeat: direct the reader to the project thesis for full details
    \item Summarize key results: validated robot model, reach/place task performance, workspace analysis findings
    \item State what is taken as given for the master thesis
\end{itemize}

\section{Cloth and Deformable Object Simulation}
\begin{itemize}
    \item Why deformable objects break rigid-body assumptions
    \item Common cloth modeling approaches: mass-spring, particle-based (PBD), continuum (FEM/MPM)
    \item Position-Based Dynamics (PBD) in PhysX: particles, constraints, Gauss-Seidel projection, XPBD compliance
    \item Vertex Block Descent (VBD) / Newton Thin Deformables: implicit energy minimization, vertex-level Newton updates
    \item Comparison: PBD vs.\ VBD for cloth in RL workflows (speed, stability, stiffness fidelity)
    \item Isaac Sim cloth integration: particle cloth API, material parameters, collision with articulations
\end{itemize}

\section{Multi-Agent Reinforcement Learning}
\begin{itemize}
    \item Extension from single-agent MDP to multi-agent: Dec-POMDP, centralized training / decentralized execution (CTDE)
    \item Multi-agent PPO variants: independent PPO (IPPO), multi-agent PPO (MAPPO)
    \item Shared vs.\ independent reward structures; cooperative vs.\ mixed objectives
    \item Communication and coordination mechanisms in multi-agent RL
    \item Prior work: multi-robot manipulation (e.g., bi-manual grasping, handover tasks)
\end{itemize}

\section{Vision-Based Classification for Robotic Sorting}
\begin{itemize}
    \item Visual inspection in manufacturing and recycling contexts
    \item CNN-based defect detection and classification (ResNet, EfficientNet architectures)
    \item Sim-to-real for visual classifiers: synthetic data generation, domain randomization for textures/lighting/backgrounds
    \item Integration of perception with robot control: observation pipelines in IsaacLab
\end{itemize}

\section{Prior Work in Robotic Cloth Manipulation}
\begin{itemize}
    \item Robotic cloth folding, unfolding, smoothing: key benchmarks (GarmentBench, GarmentLab, Seita et al.)
    \item RL approaches for deformable manipulation: reward design challenges, representation learning for cloth state
    \item Sim-to-real for cloth tasks: domain randomization of material properties, teacher-student distillation
    \item Multi-robot cloth handling: bi-manual stretching, cooperative folding
    \item Gap: no existing work combines tendon-driven robots, multi-agent RL, and condition-based sorting
\end{itemize}

\section{Sim-to-Real Transfer Strategies}
\begin{itemize}
    \item Domain randomization: which parameters and how aggressively
    \item System identification and parameter calibration
    \item Teacher-student policy distillation
    \item ROS 2 as the deployment middleware: joint state interfaces, camera topics, action servers
    \item Prior sim-to-real demonstrations for manipulation tasks
\end{itemize}

%--- ~3--4 pages ---%
\chapter{Problem Statement and Research Gap}\label{ch:research_gap}
\markboth{Problem Statement and Research Gap}{}

\section{Critical Analysis of State of the Art}
\begin{itemize}
    \item Cloth manipulation studied primarily with rigid-joint robots; tendon-driven compliance advantages unexplored for deformable objects
    \item Multi-agent RL applied to rigid object manipulation but not to cloth sorting pipelines
    \item Vision-based classification exists for quality inspection but not integrated into a multi-agent pick-stretch-sort loop
    \item Sim-to-real for multi-agent deformable manipulation is an open problem
\end{itemize}

\section{Research Gap and Justification}
\begin{itemize}
    \item Key gap: no existing system combines multi-agent RL, tendon-driven manipulation, deformable cloth physics, and condition-based classification in a single pipeline
    \item Scientific contribution: multi-agent MDP formulation for condition-based textile sorting; evaluation of cooperative policies for pick-stretch-classify-sort
    \item Practical contribution: complete simulation framework for textile recycling automation; ROS 2 integration pathway
\end{itemize}

\section{Refined Research Question}
\begin{itemize}
    \item How can a multi-agent RL system coordinate a tendon-driven picker and a collaborative sorting arm to selectively retrieve, inspect, and sort textiles by damage condition, and what strategies enable robust sim-to-real transfer?
    \item Expected contributions: multi-agent sorting pipeline, cloth sim integration, classification integration, sim-to-real preparation
\end{itemize}

%--- ~20--25 pages ---%
\chapter{Methodology}\label{ch:methodology}
\markboth{Methodology}{}

\section{Overview of Approach}
\begin{itemize}
    \item System architecture diagram: conveyor $\rightarrow$ tensegrity picker $\rightarrow$ handoff $\rightarrow$ sorting arm $\rightarrow$ stretch $\rightarrow$ camera $\rightarrow$ classify $\rightarrow$ sort into bins
    \item Development stages: (1) sorting task, (2) cloth integration, (3) second arm, (4) classification, (5) multi-agent, (6) sim-to-real prep
\end{itemize}

\section{Sorting Task (Single-Agent Extension)}
\begin{itemize}
    \item Extend cube sorting to selective retrieval: T-shirt vs.\ non-target objects on moving conveyor
    \item MDP formulation: observation (joint states, conveyor state, object labels/features), action (joint targets + gripper), reward (correct pick, correct ignore, penalty for false pick/miss)
    \item Training with domain randomization (belt speed, object placement, friction)
\end{itemize}

\section{Cloth Simulation Integration}
\begin{itemize}
    \item T-shirt cloth model: particle mesh, material parameters (mass, stretch stiffness, bending stiffness, friction, damping)
    \item PBD vs.\ VBD selection and justification for training throughput
    \item Cloth-articulation interaction: gripper grasping cloth particles, collision handling
    \item Validation of cloth behavior: qualitative draping tests, quantitative stretch tests
    \item Domain randomization of cloth parameters
\end{itemize}

\section{Second Robot Arm Integration}
\begin{itemize}
    \item Arm selection and kinematic configuration
    \item Scene placement: workspace complementarity with tensegrity picker
    \item Stretching primitive: grasp cloth edge, extend to flat configuration
    \item Camera mounting and field-of-view design
\end{itemize}

\section{Camera-Based Damage Classification}
\begin{itemize}
    \item Simulated camera setup: resolution, positioning, synthetic image generation
    \item Damage classes: intact, minor damage, major damage, unusable
    \item Classifier architecture: CNN (e.g., ResNet-18) trained on synthetic stretched-cloth images
    \item Training data generation: procedural damage textures, domain randomization (lighting, background, cloth color)
    \item Integration as observation in the RL loop or as a post-stretch evaluation module
\end{itemize}

\section{Multi-Agent RL Formulation}
\begin{itemize}
    \item Agent decomposition: Agent 1 (tensegrity picker), Agent 2 (sorting arm)
    \item Observation spaces per agent; shared state components
    \item Action spaces per agent
    \item Reward structure: cooperative (shared sorting reward), individual shaping rewards per agent
    \item Coordination protocol: sequential (pick $\rightarrow$ handoff $\rightarrow$ stretch $\rightarrow$ classify $\rightarrow$ sort) or learned timing
    \item Algorithm: MAPPO or IPPO via skrl; centralized critic with decentralized actors
    \item Curriculum: single-agent sub-tasks $\rightarrow$ coordinated full pipeline
\end{itemize}

\section{Sim-to-Real Transfer Preparation}
\begin{itemize}
    \item Domain randomization strategy: systematic parameter sweep and sensitivity analysis
    \item ROS 2 architecture: joint state publisher/subscriber, policy inference node, camera pipeline, safety constraints
    \item Action-space compatibility: simulation actions $\rightarrow$ ROS 2 joint trajectory commands
    \item Latency and timing considerations for real-time execution
    \item Gap analysis: what remains to be addressed for physical deployment
\end{itemize}

%--- ~15--20 pages ---%
\chapter{Results}\label{ch:results}
\markboth{Results}{}

\section{Sorting Task Results (Single Agent)}
\begin{itemize}
    \item Sorting accuracy: T-shirt correctly retrieved, non-targets correctly ignored
    \item False-pick rate, missed-pick rate
    \item Learning curves; comparison with project thesis cube sort baseline
    \item Effect of domain randomization on robustness
\end{itemize}

\section{Cloth Simulation Results}
\begin{itemize}
    \item Cloth behavior validation: draping, stretching, grasping fidelity
    \item Performance impact: training FPS with cloth vs.\ rigid objects
    \item Stability under domain randomization of cloth parameters
\end{itemize}

\section{Multi-Agent Sorting Pipeline Results}
\begin{itemize}
    \item End-to-end sorting accuracy by damage class
    \item Cycle time: pick to sort completion
    \item Coordination efficiency: handoff success rate, idle time
    \item Learning curves for multi-agent vs.\ single-agent baselines
    \item Qualitative policy behavior: keyframe sequences of the full pipeline
\end{itemize}

\section{Damage Classification Results}
\begin{itemize}
    \item Classification accuracy per damage class (confusion matrix)
    \item Effect of cloth stretching quality on classification accuracy
    \item Robustness under visual domain randomization
\end{itemize}

\section{Ablation Studies}
\begin{itemize}
    \item Without stretching: classification accuracy on unstretched cloth
    \item Without classification: sorting by pick heuristics only
    \item Single-agent vs.\ multi-agent: performance comparison
    \item Shared vs.\ independent reward: coordination quality
\end{itemize}

\section{Sim-to-Real Readiness Assessment}
\begin{itemize}
    \item Robustness metrics under maximal domain randomization
    \item ROS 2 interface validation in simulation (latency, message integrity)
    \item Identified gaps and risk assessment for physical deployment
\end{itemize}

%--- ~5--7 pages ---%
\chapter{Discussion}\label{ch:discussion}
\markboth{Discussion}{}

\section{Evaluation of Outcomes}
\begin{itemize}
    \item Achievement of each thesis objective
    \item Multi-agent vs.\ single-agent: when is coordination worth the complexity?
    \item Cloth simulation fidelity: sufficient for policy learning? What artifacts affect results?
\end{itemize}

\section{Comparison with Related Work}
\begin{itemize}
    \item Comparison with GarmentBench/GarmentLab benchmarks
    \item Comparison with single-robot cloth manipulation approaches
    \item Novelty: first multi-agent tendon-driven cloth sorting pipeline
\end{itemize}

\section{Limitations}
\begin{itemize}
    \item Cloth simulation approximations (PBD/VBD vs.\ real fabric behavior)
    \item Synthetic damage textures vs.\ real-world garment damage
    \item No physical deployment; sim-to-real gap remains
    \item Classification limited to synthetic training data
    \item Computational cost of multi-agent training with cloth
\end{itemize}

\section{Practical Implications}
\begin{itemize}
    \item Applicability to textile recycling industry
    \item Modularity: individual pipeline components (classifier, sorting policy) can be updated independently
    \item ROS 2 integration pathway enables incremental physical deployment
\end{itemize}

%--- ~3 pages ---%
\chapter{Conclusion}\label{ch:conclusion}
\markboth{Conclusion}{}

\section{Conclusion}
\begin{itemize}
    \item Summary of the complete multi-agent textile sorting pipeline
    \item Key quantitative achievements: sorting accuracy, classification accuracy, cycle time
    \item Answer the research question
\end{itemize}

\section{Outlook}
\begin{itemize}
    \item Physical deployment via ROS 2 on real hardware
    \item Real-world fine-tuning: sim-to-real domain adaptation, real cloth training data
    \item More garment types and damage categories
    \item Higher-level task planning for continuous conveyor operation
    \item Integration with upstream waste detection and downstream recycling processes
    \item More realistic tendon model incorporating cable dynamics
\end{itemize}
